\documentclass[11pt,a4paper]{article}
\usepackage[T1]{fontenc}
\usepackage[utf8]{inputenc}
\usepackage{lmodern,amsmath,amssymb}
\usepackage[margin=25mm]{geometry}
\usepackage[hidelinks]{hyperref}
\hypersetup{pdftitle={Topology, contraction and the missing action scale},pdfauthor={navstokgap research working notes}}
\newcommand{\tightlist}{\setlength{\itemsep}{0pt}\setlength{\parskip}{0pt}}
\setlength{\emergencystretch}{3em}
\setcounter{secnumdepth}{0}
\begin{document}
\hypertarget{a-topological-energy-floor-survives-shrinking-while-the-action-cost-vanishes}{%
\section{A topological energy floor survives shrinking while the action
cost
vanishes}\label{a-topological-energy-floor-survives-shrinking-while-the-action-cost-vanishes}}

The classical 2+1 dimensional O(3) sigma field has smooth degree-one
static configurations of every radius, all with energy \(4\pi\rho\).
Their action magnitude over one radius-crossing time is
\(4\pi\rho R/c\), which tends to zero as the radius \(R\) shrinks. Thus
a fixed topological sector can exclude zero energy without selecting a
positive action floor. The same physical field has arbitrarily
low-frequency vacuum waves.

Status: Q04 exploratory construction and written test; no
accepted-ledger promotion or novelty claim. The sector bound and scale
family are established sigma-model ingredients. The action and gap
comparisons below are the task's use of them. This continuum classical
field is a proposed physical model; it is not derived from the
repository's Newtonian particle mechanics.

\hypertarget{field-action-and-admissible-sector}{%
\subsection{1. Field, action and admissible
sector}\label{field-action-and-admissible-sector}}

Take a dimensionless unit vector \(n(t,x)\in S^2\), \(x\in\mathbb R^2\),
with propagation speed \(c>0\) and stiffness \(\rho>0\) of energy units.
Specify physical dynamics by the Lorentzian action

\[S[n]=\frac\rho2\int dt\,d^2x\,
\left(c^{-2}|\partial_tn|^2-|\nabla n|^2\right). \tag{1}\]

Here \(S\) has energy-times-time units, and the canonical momentum
density is \(\pi=\rho\partial_tn/c^2\). The Hamiltonian is

\[H=\frac\rho2\int_{\mathbb R^2}
\left(c^{-2}|\partial_tn|^2+|\nabla n|^2\right)d^2x. \tag{2}\]

Require finite energy and smooth extension of each spatial field to the
one-point compactification, with \(n(\infty)=e_3\). Its degree is

\[Q=\frac1{4\pi}\int n\cdot(\partial_1n\times\partial_2n)\,d^2x
\in\mathbb Z. \tag{3}\]

Choose the sector \(|Q|=1\). Degree is conserved during smooth
evolutions preserving the boundary extension. This is a restriction on
admitted initial data and regular evolution; it does not explain why
nature excludes \(Q=0\). No lattice spacing, lower radius or time window
is imposed.

\hypertarget{square-completion-and-an-explicit-shrinking-family}{%
\subsection{2. Square completion and an explicit shrinking
family}\label{square-completion-and-an-explicit-shrinking-family}}

Put \(a=\partial_1n\), \(b=\partial_2n\) and \(q=n\cdot(a\times b)\).
Tangency to the unit sphere gives

\[|a+n\times b|^2=|a|^2+|b|^2-2q,\qquad
|a-n\times b|^2=|a|^2+|b|^2+2q.\]

Integrating the appropriate square proves

\[H\ge E[n]:=\frac\rho2\int|\nabla n|^2d^2x
\ge4\pi\rho|Q|. \tag{4}\]

For every \(R>0\) set \(r^2=x_1^2+x_2^2\) and

\[n_R(x)=\frac{(2Rx_1,\,2Rx_2,\,r^2-R^2)}{r^2+R^2}. \tag{5}\]

Its norm is one, its boundary value is \(e_3\), and direct
differentiation gives

\[|\nabla n_R|^2=\frac{8R^2}{(r^2+R^2)^2},\qquad
q_R=-\frac{4R^2}{(r^2+R^2)^2}.\]

Using \(d^2x=2\pi r\,dr\) and
\(\int_0^\infty r(r^2+R^2)^{-2}dr=1/(2R^2)\) yields \(Q=-1\) and
\(E[n_R]=4\pi\rho\). Spatial reflection gives the opposite degree.
Equality in (4) makes each field a static energy minimizer in its
sector, hence a stationary solution of (1). This constructs actual
solutions, not only variational trial configurations.

The energy inside a disk of radius \(r\) is

\[E_R(r)=4\pi\rho\frac{r^2}{r^2+R^2}. \tag{6}\]

Consequently \(R\) is the half-energy radius, an intrinsic size rather
than an arbitrary coordinate label. In the contraction \(R\downarrow0\),
all members remain smooth and in the same sector. The limiting field is
singular at the origin: away from it \(n_R\to e_3\), while
\(n_R(0)=-e_3\). Their energy measures converge weakly to
\(4\pi\rho\delta_0\) and their gradients do not converge strongly to the
vacuum in \(L^2\). Excluding the singular endpoint does not exclude the
smooth sequence. This is a family of stationary preparations; no
dynamical collapse or finite-inertia dilation mode is asserted.

\hypertarget{which-action-is-being-tested}{%
\subsection{3. Which action is being
tested?}\label{which-action-is-being-tested}}

Define the candidate size-based action observable on this static family
by

\[\mathcal A_R:=E[n_R]\,\tau_R,\qquad
\tau_R:=R/c.\]

For the static solution on a physical time interval of length
\(\tau_R\), \(S=-E\tau_R\), so this observable equals its action
magnitude:

\[\mathcal A_R=|S[n_R;0,R/c]|=\frac{4\pi\rho R}{c}
\longrightarrow0. \tag{7}\]

A crossing time measures propagation over the half-energy radius; it is
not an oscillation period of this stationary field. Thus (7) tests this
specified energy-time candidate, not every functional with action units.
Equivalently, under spacetime dilation
\(n_\lambda(t,x)=n(t/\lambda,x/\lambda)\) on a time interval scaled by
\(\lambda\), (1) gives \(S[n_\lambda]=\lambda S[n]\), whereas spatial
degree is unchanged. The field equation retains this dilation symmetry
at fixed \(\rho,c\).

For a fixed externally prescribed duration \(T>0\), the same static
configurations, each held for \(T\), have \(|S|=ET\ge4\pi\rho T\). Its
units and lower bound include that supplied clock. A canonical motion
cost \(\int dt\int\pi\cdot\partial_tn\,d^2x\) is exactly zero on every
member of (5), even though its spatial gradient energy is positive.
Topology excludes a spatially uniform field, but admits rest in
configuration space.

The available dimensional parameters \(\rho\) and \(c\) cannot by
themselves form an action: a length or a time must also enter. A lower
radius \(R\ge\ell>0\) would give \(\mathcal A_R\ge4\pi\rho\ell/c\)
within this family. That adds an independently justified physical scale
and its coefficient; its survival as \(\ell\downarrow0\) would need a
new mechanism. No physical-time attraction, preparation-independent
constant or quantum phase rule follows from the stationary family.

\hypertarget{the-sector-energy-floor-is-not-a-frequency-gap}{%
\subsection{4. The sector energy floor is not a frequency
gap}\label{the-sector-energy-floor-is-not-a-frequency-gap}}

The physical action (1), rather than an auxiliary sampling generator,
determines small fluctuations. Near the uniform vacuum write
\(n=(u_1,u_2,\sqrt{1-|u|^2})\). The quadratic action gives

\[\partial_t^2u=c^2\Delta u,\qquad \omega(k)=c|k|. \tag{8}\]

The squared-frequency operator \(-c^2\Delta\) on
\(L^2(\mathbb R^2;\mathbb R^2)\) with domain \(H^2\) has spectrum
\([0,\infty)\). For example normalized smooth packets
\(u_L(x)=L^{-1}f(x/L)\) have fixed \(L^2\) norm and Rayleigh quotient
\(c^2L^{-2}\|\nabla f\|_2^2/\|f\|_2^2\to0\). On a periodic box of side
\(L\), after removing constant rotations, the first frequency is
\(2\pi c/L\); the thermodynamic limit closes it. These are vacuum
linearized statements, not a quantized Hamiltonian spectrum.

Even within the degree-one sector there is no isolated static minimizer:
translations \(n_R(x-a)\) retain the same energy. Their tangent vectors
\(\partial_i n_R\) are square integrable by (4) and are zero directions
of the static second variation. Quotienting translations is a further
restriction; no positive internal fluctuation gap after such a quotient
is claimed here. In particular (4) compares a chosen sector to the
vacuum, not all excitations to a unique ground state within one sector.

\hypertarget{decision-and-next-physical-construction}{%
\subsection{5. Decision and next physical
construction}\label{decision-and-next-physical-construction}}

Retain topology as a conservative mechanism for a positive \emph{sector
energy} at fixed stiffness. Park topology alone as an action selector in
this model: it leaves an arbitrary radius, an admitted static state and
low-frequency vacuum waves. This advances the zero-exclusion test beyond
Q03's need for continuous power, while exposing separate size and motion
requirements.

A concrete next candidate is a unit-vector field in three spatial
dimensions with both quadratic and quartic spatial-gradient energies.
Under a size rescaling their contributions behave as \(aR\) and \(b/R\).
Their competition can fix a radius without external fueling. The next
decision is whether this stabilization supplies an action or only
expresses a scale already in \(a,b\) and the kinetic normalization; it
must also test the vacuum spectrum. Existence of a minimizer cannot be
inferred from the one-parameter balance alone. This is the selected
candidate test, not a new theorem or a claim that stabilization
establishes quantum necessity.

\hypertarget{source-and-review}{%
\subsection{Source and review}\label{source-and-review}}

M. S. Ody and L. H. Ryder, \emph{Time-Independent Solutions to the
Two-Dimensional Non-Linear O(3) Sigma Model and Surfaces of Constant
Mean Curvature},
\href{https://arxiv.org/abs/hep-th/9402137v1}{arXiv:hep-th/9402137v1}
(1994), section 2, pp.~4--7, equations (2.1)--(2.26), supplies the
static degree, energy and arbitrary-scale ingredients. Its unit
stiffness gives \(E=4\pi|Q|\) for the duality solutions. The Lorentzian
dynamics, the radius-crossing observable and the physical gap comparison
are specified here.

The \href{../reviews/topological-sector-Q04.md}{Q04 review} records the
written proof check and bounded primary-source comparison: one
Sol-medium worker, two queries, one primary paper and six pages, with
formula-image checks. The coordinator rechecked the square completion,
radial integral, dilation and wave normalization, and visually checked
source pp.~6--7. No numerical or symbolic verification scripts were
used. Source matches concern the static ingredients; no novelty or
accepted-ledger promotion is asserted.
\end{document}
